% Generated draft: requires mathematical and source review.
\documentclass[11pt]{article}
\usepackage[margin=25mm]{geometry}
\usepackage{fontspec}
\setmainfont{Latin Modern Roman}
\usepackage{amsmath,amssymb,mathtools,unicode-math}
\setmathfont{Latin Modern Math}
\usepackage{enumitem,needspace,xurl,hyperref}
\setlength{\parindent}{0pt}
\setlength{\parskip}{5pt}
\setlength{\emergencystretch}{3em}
\newcommand{\N}{\mathbb{N}}
\newcommand{\Z}{\mathbb{Z}}
\newcommand{\Q}{\mathbb{Q}}
\newcommand{\R}{\mathbb{R}}
\newcommand{\C}{\mathbb{C}}
\newcommand{\F}{\mathbb{F}}
\newcommand{\A}{\mathbb{A}}
\newcommand{\PP}{\mathbb{P}}
\newcommand{\E}{\mathbb{E}}
\newcommand{\eps}{\varepsilon}
\newcommand{\dd}{\,\mathrm{d}}
\begin{document}
\section*{UnsolvedMath ID: 701224}
\textbf{Source problem code:} catalogue entry
\subsection{Definitions and mathematical statement}
A finite composition series of a group \(G\) is a chain
\[
1=G_0\triangleleft G_1\triangleleft\cdots\triangleleft G_n=G,
\]
where \(n\geq 0\) and every factor \(G_i/G_{i-1}\), for \(1\leq i\leq n\), is a nontrivial simple group. The existence of such a series does not require \(G\) or its composition factors to be finite.

The source does not specify the ambient universe or closure conventions for Fitting classes. The following makes the question explicit relative to an ambient class \(\mathcal U\) of groups and the finite normal-product convention. A Fitting class in \(\mathcal U\) is a nonempty subclass \(\mathcal F\subseteq\mathcal U\), closed under isomorphism, such that:
\begin{enumerate}
\item if \(H\in\mathcal F\) and \(N\triangleleft H\), then \(N\in\mathcal F\);
\item if \(N_1,N_2\triangleleft H\), \(H\in\mathcal U\), \(H=N_1N_2\), and \(N_1,N_2\in\mathcal F\), then \(H\in\mathcal F\).
\end{enumerate}
Assume that \(\mathcal U\) is closed under isomorphism and normal subgroups, so that it is itself a Fitting class under these conventions. For \(G\in\mathcal U\), define the Fitting class generated by \(G\) by
\[
\operatorname{Fit}_{\mathcal U}(G)
=
\bigcap\{\mathcal F\subseteq\mathcal U:
\mathcal F\text{ is a Fitting class and }G\in\mathcal F\}.
\]
A Fitting subclass of \(\operatorname{Fit}_{\mathcal U}(G)\) means a Fitting class \(\mathcal H\) in the same ambient universe satisfying \(\mathcal H\subseteq\operatorname{Fit}_{\mathcal U}(G)\); equality is allowed.

\textbf{Question (V. A. Vedernikov, Kourovka Notebook 16.32).} For every group \(G\) having a composition series, is the collection
\[
\{\mathcal H\subseteq\operatorname{Fit}_{\mathcal U}(G):
\mathcal H\text{ is a Fitting class in }\mathcal U\}
\]
finite? Here \(G\) must belong to the intended ambient universe \(\mathcal U\). This is a conditional rendering of the supplied question: the source writes \(\operatorname{Fit}(G)\) without specifying \(\mathcal U\) or whether additional closure under infinite normal products is intended. Those conventions require confirmation before this rendering can be treated as definitive. The source also says to compare Problem 14.31, whose statement was not supplied.
\subsection{Short English statement}
Vedernikov asks whether a group with a finite composition series generates a Fitting class containing only finitely many Fitting subclasses. The group need not have finite order. The ambient universe and the intended Fitting-class closure conventions remain to be confirmed.
\subsection{Sources}
\begin{itemize}
\item Local UnsolvedMath registry, supplied complete source record: permanent ID \(701224\), exact source code \texttt{catalogue entry}, title \emph{Kourovka Notebook Problem 16.32}; upstream revision \texttt{d371ee917ac9647fbb06cd6fc529deec6f3dda82}. The supplied registry labels the problem \emph{open}; this label is not an independently verified mathematical or literature assessment.
\item Supplied source attribution: \emph{The Kourovka Notebook (Unsolved Problems in Group Theory)}, 16th issue (2006), Problem 16.32, proposed by V. A. Vedernikov. The record also identifies the 21st issue edition (2026), \texttt{arXiv:1401.0300v48}, and names E. I. Khukhro and V. D. Mazurov as editors of the notebook collection. This bibliographic attribution is supplied evidence and has not been independently checked.
\item The statement includes the cross-reference ``Cf. 14.31.'' No text of Problem 14.31 or associated explanatory material was supplied.
\end{itemize}
\end{document}
